\section*{Positive $U$ tax absorbs the growing-rank OR boundary layer}

All logarithms are base two.  For integers $r\ge2$ and
$k\ge1$, split the cube into $k$ disjoint $r$-bit blocks.  Let $D_j$
be the event that block $j$ equals a fixed point and put
\[
 f_{k,r}=1-2\mathbf 1_{D_1\vee\cdots\vee D_k}.
 \tag{1}
\]
This is a genuine Boolean family, not a tensor power.  Fix
$\lambda>0$, write $\epsilon=2^{-r}$, and take
$k_r=\lfloor\lambda/\epsilon\rfloor$.  The frozen GU pair in this
packet is
\[
 (a,b)=\left({3\over4},{1\over4}\right).
 \tag{2}
\]

\paragraph{Theorem.}
For every fixed $\lambda>0$,
\[
 \begin{aligned}
 I(f_{k_r,r})&=2\lambda e^{-\lambda}r+O_\lambda(r2^{-r}),\\
 H(f_{k_r,r})&=8\lambda e^{-\lambda}r+O_\lambda(1),\\
 U(f_{k_r,r})&=2\lambda e^{-\lambda}r+O_\lambda(1),\\
 G(f_{k_r,r})&=2\lambda e^{-\lambda}r+O_\lambda(r^2 2^{-r}).
 \end{aligned}                                                   \tag{3}
\]
More precisely,
\[
 U-I\longrightarrow
 \Psi(\lambda):=
 c\bigl((2e^{-\lambda}-1)^2\bigr)
 +{4\lambda e^{-\lambda}(1-e^{-\lambda})\over\ln2}>0,           \tag{4}
\]
where $c(x)=x\log_2(1/x)$ and $c(0)=0$, while
\[
 G-I\longrightarrow0.                                            \tag{5}
\]
Consequently, for every fixed $b\in(0,1)$,
\[
 \boxed{
 G-(1-b)I-bU\longrightarrow-b\Psi(\lambda)<0.}                  \tag{6}
\]
In particular, the frozen residual satisfies
\[
 \boxed{
 G-{3\over4}I-{1\over4}U
 \longrightarrow-{1\over4}\Psi(\lambda)<0.}                    \tag{7}
\]
Thus for each fixed $\lambda>0$ there exists
$r_0(\lambda,b)$ such that the boundary pair $(1-b,b)$ holds on this
family for all $r\ge r_0$.  The threshold is not asserted uniform as
$b\downarrow0$ or as $\lambda$ approaches $0$ or infinity.

\paragraph{Exact formulas.}
Put $P=(1-\epsilon)^k$, $p_0=(2P-1)^2$, and
$L=-\log_2(1-\epsilon)$.  Exact Fourier and upper-set summation gives
\[
 I={2Pkr\epsilon\over1-\epsilon},                                 \tag{8}
\]
\[
 \begin{aligned}
 H={}&c(p_0)-2(1-p_0)+8Pkr\epsilon
       +8PkL\bigl((1-\epsilon)-P\bigr),\\
 U={}&c(p_0)+4Pk\epsilon
       \left(r-{r\over2(1-\epsilon)}\right)
       +4PkL\bigl((1-\epsilon)-P\bigr).
 \end{aligned}                                                    \tag{9}
\]
Restriction of one physical bit has unequal children: one deletes its
block, while the other replaces it by an $(r-1)$-bit defect.  If
$U_{k,r}$ denotes (9), then exactly
\[
 G=kr\left(U_{k,r}
 -{U_{k-1,r}+U_{(r-1,r,\ldots,r)}\over2}\right).                 \tag{10}
\]
Thus every normalization and terminal term, and both restriction
children, remain in the calculation.

\paragraph{First-order and terminal limits.}
Since $P=e^{-\lambda}(1+O_\lambda(\epsilon))$, (8)--(9) yield (3).
Subtracting (8) from the expression for $U$ and taking the limit gives
exactly (4).  Its second summand is strictly positive for every
$\lambda>0$, including the balanced value $\lambda=\ln2$ at which
$p_0\to0$.

It remains to justify (5), since the ratio statement $G/I\to1$ alone
would not suffice.  Away from $\lambda=\ln2$, the constant-atom entropy
is smooth in a fixed neighbourhood of $P=e^{-\lambda}$.  Expanding the
two child products
\[
 P_-={P\over1-\epsilon},\qquad
 P_+={P(1-2\epsilon)\over1-\epsilon},\qquad
 {P_-+P_+\over2}=P                                      \tag{11}
\]
in (9)--(10) gives
\[
 {G\over I}=1+2r\epsilon+O_\lambda(\epsilon),
 \qquad
 G-I=4\lambda e^{-\lambda}r^2\epsilon+O_\lambda(r\epsilon).
 \tag{12}
\]
Both terms in the last display tend to zero.

At $\lambda=\ln2$, $c((2P-1)^2)$ is not smooth and must not be put into
the error in (12).  Write
$k=\lambda/\epsilon+\alpha_r$, $-1<\alpha_r\le0$, and
$q_r=\lambda/2+\alpha_r$.  Then the parent, deleted-child, and
shortened-child constant Fourier amplitudes are respectively
\[
 -q_r\epsilon+O(\epsilon^2),\quad
 -(q_r-1)\epsilon+O(\epsilon^2),\quad
 -(q_r+1)\epsilon+O(\epsilon^2).
\]
Their central entropy difference equals
\[
 c(q_r^2\epsilon^2)
 -{c((q_r-1)^2\epsilon^2)+c((q_r+1)^2\epsilon^2)\over2}
 =-2r\epsilon^2+O(\epsilon^2),                         \tag{13}
\]
uniformly in the bounded $q_r$.  After multiplication by $kr$, this
cancels the $+2r\epsilon$ normalized correction from the nonconstant
rows.  Hence
\[
 {G\over I}=1+O(\epsilon),\qquad G-I=O(r\epsilon)\to0,             \tag{14}
\]
which proves (5) at the exceptional intensity as well.

\paragraph{GU implication and limitation.}
The same expansion shows $G>I$ eventually whenever
$\lambda\ne\ln2$, so the endpoint $(a,b)=(1,0)$ is false on this
family.  Equations (4)--(7) show that this endpoint failure is only a
vanishing boundary layer and cannot refute any fixed boundary pair with
$b>0$.  No ratio diverges: $G/I\to1$, $U/I\to1$, and $H/I\to4$.
Therefore this packet removes a proposed uniformity obstruction but does
not prove universal GU or FEI.  What remains is to obtain one fixed
positive fraction of the terminal reserve for arbitrary Boolean
functions without paying the same influence more than once.
